\section{Prompt Templates}\label{appendix-d-prompt-templates}

This appendix documents the exact prompt templates used for waypoint elicitation across the benchmark's 11 phases.

\subsection{D.1 Standard Elicitation Prompts}\label{d.1-standard-elicitation-prompts}

Two prompt formats were developed and evaluated during Phase 1 prompt selection. The semantic format was selected as the primary template based on higher extraction accuracy and more interpretable waypoints on the holdout set.

\subsubsection{Semantic Format (Primary)}\label{semantic-format-primary}

\begin{verbatim}
Imagine walking through conceptual space from "{A}" to "{B}".
What {N} landmarks do you pass along the way?
List only the landmark concepts, one per line, numbered 1 through {N}.
Do not include the starting or ending concepts.
\end{verbatim}

\subsubsection{Direct Format (Fallback)}\label{direct-format-fallback}

\begin{verbatim}
List exactly {N} intermediate concepts that form a path from "{A}" to "{B}".
Respond with only the concepts, one per line, numbered 1 through {N}.
Do not include the starting or ending concepts.
\end{verbatim}

The direct format was used as a fallback for models that produced unparseable outputs under the semantic prompt, notably during Phase 10A reliability probing of new models.

\textbf{Default parameters.} Unless otherwise noted: N = 7 waypoints, temperature = 0.7, 10--20 independent runs per (model, pair, direction) condition. Phases 1--2 used 20 runs; later phases standardized to 10 runs.

\subsection{D.2 Pre-Fill Elicitation Prompt}\label{d.2-pre-fill-elicitation-prompt}

Introduced in Phase 7A for causal intervention experiments. The first waypoint is pre-specified, and the model generates the remaining N-1 waypoints.

\begin{verbatim}
The first intermediate concept between "{A}" and "{B}" is "{pre-fill}".
List exactly {N-1} more intermediate concepts that continue the path
from "{pre-fill}" to "{B}".
Respond with only the concepts, one per line, numbered 1 through {N-1}.
Do not include "{A}", "{pre-fill}", or "{B}" in your list.
\end{verbatim}

With default parameters, N-1 = 6 (the pre-fill occupies position 1, and the model generates positions 2--7).

\subsection{D.3 System Prompt}\label{d.3-system-prompt}

All elicitations used the following system prompt:

\begin{verbatim}
You are a helpful assistant that provides single-word or short-phrase
responses representing conceptual landmarks.
\end{verbatim}

This system prompt was consistent across all models and phases. No chain-of-thought or reasoning instructions were included; the waypoint task is presented as a direct generation task.

\subsection{D.4 API Configuration}\label{d.4-api-configuration}

All models were accessed via the OpenRouter API (\texttt{https://openrouter.ai/api/v1/chat/completions}) with the following parameters:

\begin{itemize}
\tightlist
\item
  \texttt{max\_tokens}: 256
\item
  \texttt{temperature}: 0.7 (default; varied in Phase 11C robustness analysis)
\item
  \texttt{top\_p}: 1.0
\item
  \texttt{frequency\_penalty}: 0
\item
  \texttt{presence\_penalty}: 0
\end{itemize}

Phase 10A introduced a \texttt{-\/-patient} CLI flag that extended the request timeout from 60 seconds to 300 seconds, recovering 3 of 5 new models that failed the default timeout on OpenRouter latency.

\subsection{D.5 Response Parsing}\label{d.5-response-parsing}

Model responses were parsed using a multi-strategy extraction pipeline attempting, in order:

\begin{enumerate}
\def\labelenumi{\arabic{enumi}.}
\tightlist
\item
  JSON array extraction
\item
  Numbered list parsing (e.g., ``1. concept'')
\item
  Bullet list parsing (e.g., ``- concept'')
\item
  Comma-separated value extraction
\item
  Arrow-separated extraction (recognizing \texttt{-\textgreater{}}, \texttt{=\textgreater{}}, \texttt{-\/-\textgreater{}}, \texttt{==\textgreater{}})
\item
  Line-by-line fallback
\end{enumerate}

After parsing, responses passed through the canonicalization pipeline (Section 2.3 of the main text): lowercase, article stripping, lemmatization via the compromise NLP library, and phrase normalization.

\subsection{D.6 Prompt Variations Across Phases}\label{d.6-prompt-variations-across-phases}

\begin{longtable}[]{@{}
  >{\raggedright\arraybackslash}p{(\columnwidth - 4\tabcolsep) * \real{0.2692}}
  >{\raggedright\arraybackslash}p{(\columnwidth - 4\tabcolsep) * \real{0.4231}}
  >{\raggedright\arraybackslash}p{(\columnwidth - 4\tabcolsep) * \real{0.3077}}@{}}
\toprule\noalign{}
\begin{minipage}[b]{\linewidth}\raggedright
Phase
\end{minipage} & \begin{minipage}[b]{\linewidth}\raggedright
Variation
\end{minipage} & \begin{minipage}[b]{\linewidth}\raggedright
Reason
\end{minipage} \\
\midrule\noalign{}
\endhead
\bottomrule\noalign{}
\endlastfoot
1 & Both formats tested on holdout pairs & Prompt format selection \\
1--6 & Semantic format only & Standard elicitation \\
7A--10B & Pre-fill variant added & Causal intervention \\
10A & Direct format fallback for some models & Parse reliability \\
11C & N varied (5, 7, 9); temperature varied (0.5, 0.7, 0.9) & Robustness analysis \\
\end{longtable}

No other prompt variations were introduced. The consistency of the prompt template across 11 phases is a deliberate design choice supporting cross-phase comparability.
